\documentclass[11pt]{article}

\usepackage[margin=0.85in]{geometry}
\usepackage[T1]{fontenc}
\usepackage[utf8]{inputenc}
\usepackage{lmodern}
\usepackage{hyperref}
\usepackage{longtable}
\usepackage{array}
\usepackage{enumitem}
\usepackage{parskip}
\usepackage{xcolor}
\usepackage{titlesec}

\definecolor{cpastblue}{HTML}{1D4ED8}
\definecolor{cpastgold}{HTML}{D97706}
\definecolor{cpastink}{HTML}{1F2937}
\definecolor{cpastsoft}{HTML}{EFF6FF}

\hypersetup{
  colorlinks=true,
  urlcolor=cpastblue,
  linkcolor=cpastblue
}

\titleformat{\section}{\large\bfseries\color{cpastink}}{\thesection}{0.7em}{}
\titleformat{\subsection}{\normalsize\bfseries\color{cpastink}}{\thesubsection}{0.7em}{}

\newcommand{\doclink}[2]{\href{run:#1}{#2}}
\newcommand{\weblink}[2]{\href{#1}{#2}}

\begin{document}

\begin{center}
{\LARGE\bfseries CPAST Three-Way Conference Prep Packet}\\[0.4em]
{\large Pine Brook Elementary --- Spring 2026}\\[0.3em]
\colorbox{cpastsoft}{\parbox{0.92\linewidth}{\centering
\textbf{Purpose:} Help Piter walk into the CPAST meeting ready with the right evidence, talking points, and clickable backup documents.
}}
\end{center}

\section{What The Meeting Is}

This is your \textbf{final three-way conference}. Based on your placement documents, the meeting is meant to:
\begin{itemize}[leftmargin=*]
  \item review your teaching using the CPAST rubric,
  \item compare \textbf{your self-assessment}, \textbf{Derek's assessment}, and \textbf{Zenon's assessment},
  \item discuss evidence behind ratings,
  \item work toward \textbf{consensus scores}, and
  \item identify \textbf{growth goals} for your next stage as a teacher.
\end{itemize}

\textbf{Bottom line:} it is conversational, but it is also an evidence-based evaluation meeting. Walk in ready to talk through examples, not just impressions.

\section{What To Bring Or Have Open}

\begin{itemize}[leftmargin=*]
  \item Your CPAST self-assessment notes
  \item 2--4 strong lesson examples that show planning, inclusion, assessment, and adjustment
  \item A few concrete artifacts: lesson plans, student-facing slides, reflections, screenshots, or de-identified student work
  \item One short summary of your growth this semester
  \item One short summary of what you still want to improve
\end{itemize}

\subsection{Best Clickable References}

\begin{longtable}{>{\raggedright\arraybackslash}p{0.29\linewidth} >{\raggedright\arraybackslash}p{0.67\linewidth}}
\textbf{Need} & \textbf{Open This} \\
\hline
CPAST rubric & \doclink{../warner-school-cpast-rubric.docx}{Warner School CPAST rubric} \\
Examples of evidence & \doclink{../warner-cpast-examples-of-evidence.docx}{Warner CPAST examples of evidence} \\
Field experience evaluation context & \doclink{../field-experience-eval.docx}{Field experience evaluation} \\
Placement expectations & \doclink{../expectations/letter-of-expectations-spring-2026-work.md}{Letter of Expectations work file} \\
Main action checklist & \doclink{./action-items.md}{Action items} \\
Internship summary & \doclink{./student-teaching-internship-action-items-and-key-information.md}{Student teaching internship action items and key information} \\
Placement handbook notes & \doclink{../handbook/handbook-notes.md}{Handbook notes} \\
School schedule / milestones & \doclink{../../pine-brook-elementary/spring-2026-teaching-schedule-derek.md}{Spring 2026 teaching schedule} \\
Lesson hub & \doclink{../../pine-brook-elementary/lessons/lessons-notes.md}{Lessons notes} \\
Transcript hub & \doclink{../../transcripts/transcripts-notes.md}{Transcripts notes} \\
\end{longtable}

\section{What They Are Likely Evaluating}

Your own documents repeatedly point to these evidence areas:
\begin{itemize}[leftmargin=*]
  \item planning and alignment,
  \item inclusive instruction and differentiation,
  \item formative assessment and instructional adjustment,
  \item use of digital tools in CS learning,
  \item reflective practice and professional behavior.
\end{itemize}

\subsection{Fast Evidence Checklist}

\begin{itemize}[leftmargin=*]
  \item \textbf{Planning:} a lesson with clear targets, sequence, materials, and assessment
  \item \textbf{Inclusion:} a lesson where scaffolds, visuals, chunking, or accommodations are visible
  \item \textbf{Adjustment:} an example of how you changed instruction based on what students needed
  \item \textbf{Technology/CS:} a lesson where digital tools actually supported the learning goal
  \item \textbf{Reflection:} a short example of what you learned and what you changed next
  \item \textbf{Professionalism:} consistent attendance, communication, deadlines, and responsiveness to feedback
\end{itemize}

\section{How To Talk In The Meeting}

\subsection{Good Opening}

\begin{quote}
I see this semester as growth in planning, inclusive support, and making CS learning more visible and accessible for students. I came in with one level of readiness, and I can now point to specific lessons and feedback that show how my teaching changed.
\end{quote}

\subsection{Three Strong Themes To Be Ready To Explain}

\begin{enumerate}[leftmargin=*]
  \item \textbf{How your planning improved.}
  Mention clearer targets, stronger sequencing, stronger lesson materials, or tighter alignment.
  \item \textbf{How you supported access and inclusion.}
  Mention chunking, visuals, partner structures, scaffolds, or class-specific supports.
  \item \textbf{How you responded to evidence.}
  Mention formative checks, confusion you noticed, and what you adjusted.
\end{enumerate}

\subsection{Good Reflection Formula}

Use this pattern when answering questions:
\begin{quote}
\textbf{What I was trying to teach} $\rightarrow$ \textbf{what I noticed} $\rightarrow$ \textbf{what I changed or would change next}
\end{quote}

\section{Suggested Lesson Evidence To Pull Up}

Use 2--4 lessons, not everything. Pick ones that let you talk about different strengths.

\begin{longtable}{>{\raggedright\arraybackslash}p{0.24\linewidth} >{\raggedright\arraybackslash}p{0.71\linewidth}}
\textbf{Theme} & \textbf{Good Place To Look} \\
\hline
Inclusive scaffolding & Grade-specific lesson folders under \doclink{../../pine-brook-elementary/lessons/lessons-notes.md}{Lessons notes} \\
Audience/class support needs & Class audience profiles inside each grade folder, starting from \doclink{../../pine-brook-elementary/lessons/lessons-notes.md}{Lessons notes} \\
Transcript-backed evidence & \doclink{../../transcripts/transcripts-notes.md}{Transcripts notes} \\
School expectations / curriculum spine & \doclink{../../pine-brook-elementary/teaching-placement-ignite-curriculum-gcsd-25-26/teaching-placement-ignite-curriculum-gcsd-25-26-notes.md}{IGNITE curriculum notes} \\
Schedule / chronology & \doclink{../../pine-brook-elementary/spring-2026-teaching-schedule-derek.md}{Spring 2026 teaching schedule} \\
\end{longtable}

\section{Questions They May Ask}

\begin{itemize}[leftmargin=*]
  \item Where do you see the most growth in your teaching this semester?
  \item How did you differentiate or make lessons more accessible?
  \item What evidence do you have that students understood the lesson?
  \item How did you adjust when something was not working?
  \item What area do you still want to strengthen?
\end{itemize}

\subsection{Strong Answer Strategy}

\begin{itemize}[leftmargin=*]
  \item answer briefly first,
  \item name one real lesson or class context,
  \item point to one concrete artifact if needed,
  \item end with what you learned or what you would improve.
\end{itemize}

\section{What To Avoid}

\begin{itemize}[leftmargin=*]
  \item Do not answer only in vague generalities.
  \item Do not bring too many examples; bring a small strong set.
  \item Do not treat the meeting like it is only about scores.
  \item Do not hide growth areas; name them clearly and pair them with a next step.
\end{itemize}

\section{Five-Minute Prep Before The Meeting}

\begin{enumerate}[leftmargin=*]
  \item Open the CPAST rubric.
  \item Open 2--4 lesson folders you may cite.
  \item Open one transcript or evidence note if a lesson needs backup.
  \item Write 3 bullets:
  \begin{itemize}[leftmargin=*]
    \item biggest growth,
    \item strongest evidence,
    \item next growth goal.
  \end{itemize}
  \item Keep this PDF open as your talking guide.
\end{enumerate}

\section{External Links}

\begin{itemize}[leftmargin=*]
  \item \weblink{https://osu.az1.qualtrics.com/jfe/form/SV_cC8iDwgJ6JKLl7D}{CPAST training link referenced in your task notes}
  \item \weblink{https://www.warner.rochester.edu/sites/default/files/2025-09/StudentTeachingHandbook\%20-\%202025-26.pdf}{Student Teaching Handbook}
\end{itemize}

\vfill
\begin{center}
\color{cpastgold}\rule{0.7\linewidth}{0.8pt}\\
\textbf{Meeting mindset:} calm, specific, evidence-based, reflective
\end{center}

\end{document}
